%%%%%%%%%%%%%%%%%%%%%%%%%%%%%%%%%%%%%%%%%%%%%%%%%%%%%%%%%%%%%%%%%%%%%%%%%%%%%%%%%%
\begin{frame}[fragile]\frametitle{}
\begin{center}
{\Large PCA + K-Means --- Handwritten Digit Clustering}

{\tiny (Ref: scikit-learn Decomposition and Clustering User Guide)}
\end{center}
\end{frame}

%%%%%%%%%%%%%%%%%%%%%%%%%%%%%%%%%%%%%%%%%%%%%%%%%%%%%%%%%%
\begin{frame}[fragile]\frametitle{Problem Info}
\begin{itemize}
\item Dataset: Digits (\texttt{sklearn.datasets.load\_digits}, built-in) --- 1,797 images, 64 features each
\item Part A: Use PCA to reduce 64 dimensions to 2; visualise digit clusters in 2D
\item Part B: Apply K-Means (K=10) on PCA-reduced data; measure cluster quality
\item Part C: Vary the number of PCA components (2, 10, 30) and observe effect on clustering
\end{itemize}
\end{frame}

%%%%%%%%%%%%%%%%%%%%%%%%%%%%%%%%%%%%%%%%%%%%%%%%%%%%%%%%%%
\begin{frame}[fragile]\frametitle{Import Libraries}
\begin{lstlisting}
import numpy as np
import matplotlib.pyplot as plt
from sklearn.datasets import load_digits
from sklearn.preprocessing import StandardScaler
from sklearn.decomposition import PCA
from sklearn.cluster import KMeans
from sklearn.metrics import silhouette_score, adjusted_rand_score
\end{lstlisting}
\end{frame}

%%%%%%%%%%%%%%%%%%%%%%%%%%%%%%%%%%%%%%%%%%%%%%%%%%%%%%%%%%
\begin{frame}[fragile]\frametitle{Part A: PCA Visualisation}
\begin{lstlisting}
digits = load_digits()
X = StandardScaler().fit_transform(digits.data)

pca2 = PCA(n_components=2)
X2   = pca2.fit_transform(X)

plt.figure(figsize=(8, 6))
scatter = plt.scatter(X2[:, 0], X2[:, 1],
                      c=digits.target, cmap='tab10',
                      s=10, alpha=0.7)
plt.colorbar(scatter, label='Digit')
plt.title("Digits projected to 2 PCA components")
plt.xlabel("PC 1"); plt.ylabel("PC 2")
plt.tight_layout(); plt.show()

print("Variance explained:", pca2.explained_variance_ratio_.sum())
\end{lstlisting}
\end{frame}

%%%%%%%%%%%%%%%%%%%%%%%%%%%%%%%%%%%%%%%%%%%%%%%%%%%%%%%%%%
\begin{frame}[fragile]\frametitle{Part B: K-Means on PCA Data}
\begin{itemize}
\item \textbf{Silhouette score}: measures cluster compactness (higher is better, max 1.0)
\item \textbf{Adjusted Rand Index}: compares clusters to true digit labels (1.0 = perfect)
\end{itemize}
\begin{lstlisting}
for n_comp in [2, 10, 30]:
    Xr = PCA(n_components=n_comp).fit_transform(X)
    km = KMeans(n_clusters=10, random_state=42, n_init=10)
    labels = km.fit_predict(Xr)
    sil = silhouette_score(Xr, labels)
    ari = adjusted_rand_score(digits.target, labels)
    print(f"PCA={n_comp:2d}  Silhouette={sil:.3f}  "
          f"ARI={ari:.3f}")
\end{lstlisting}
\end{frame}

%%%%%%%%%%%%%%%%%%%%%%%%%%%%%%%%%%%%%%%%%%%%%%%%%%%%%%%%%%
\begin{frame}[fragile]\frametitle{Assignment Tasks}
\begin{itemize}
\item Report the variance explained by 2, 10, and 30 PCA components
\item Plot the explained variance ratio (scree plot) and identify the ``elbow''
\item At which number of PCA components does ARI peak? Explain why
\item Visualise cluster centres (inverse-transform to pixel space) for the best configuration
\item What do the cluster centres look like? Do they resemble recognisable digits?
\end{itemize}
\end{frame}
